% Folder 83, batch 07 — archivists' pages 121–140.
% Pass: Opus 5.5 (claude-opus-5-5), 2026-10-10 — first pass, unchecked against the pages by a human.
\documentclass[11pt,a4paper]{article}
\input{../preamble/grothendieck}

\folder{83}
\batch{7}
\pages{121}{140}
\foldertitle{[Icosaèdres, bi-icosaèdres et algèbres d'Azumaya de rang 4] : notes manuscrites (1982-1983, 1986, s.d.).}
\dating{1982-1986}
\watermark{Édition de démonstration}

\begin{document}

\page{121}
\note{la page continue un argument commencé avant le lot : $\delta,\delta'$, $Q$, $\rho_0$, $b,b',c,c'$ et $t$ sont introduits plus haut.}
\struck{\uncertain{Quid P}} \ill{} deux sous-alg. \uncertain{étales} (\struck{\uncertain{diagonales}} \uncertain{stables}) $\overline{H}$, $\overline{H'}$ correspondant à l'inv. $\rho_0$. Quitte à faire une ext. \uncertain{quadratique} (\ill{}), on \supplied{peut supposer} $\overline{H}\simeq H$, $\overline{H'}\simeq H'$, \struck{\ill{}}
\marginal{OPS \struck{$H = \overline{H}$, $H'$} $H$, $H'\subset A$.}
\struck{\ill{}} Soit donc
\[ \tau = \operatorname{Tr} uv , \]
si on avait $t=\tau$, \ill{} on aurait \uncertain{fini}.

On sait que
\[ \text{\struck{$b^2-4c$}}\ \frac{\delta\delta'}{Q(t)} = \frac{\delta\delta'}{Q(\tau)} \qquad (=\rho_0) \]
donc
\[ Q(t) = Q(\tau) \qquad (=\sigma) \]
\struck{On \ill{} \uncertain{encore}, \uncertain{pour} \ill{} $k$ \uncertain{corps} alg. clos \ill{}}
\[ R(T) = T^2 - bb'T + r \]
où $r = (b^2c'+b'^2c-4cc'-\sigma)$, $t$ et $\tau$ sont des solutions de l'équation $R(T)=0$. Si donc on avait $\tau\neq t$, on aurait
\[ t+\tau = bb' \qquad\text{donc}\qquad \tau = bb'-t . \]
On considère
\[ u' = \operatorname{Tr} u.1 - u = b - u \]

\page{122}
On aura
\[ \operatorname{Tr} u' = \operatorname{Tr} u = b , \qquad \det u' = \det u = c \]
\[ \operatorname{Tr} u'v = \operatorname{Tr}(b-u)v = b\underbrace{\operatorname{Tr} v}_{b'} - \underbrace{\operatorname{Tr} uv}_{t} = bb' - t = \tau \]
et on \uncertain{gagne}, en remplaçant $u$ par $u'$.

On a donc gagné dans le cas où on est sur un corps. On \ill{} \uncertain{maintenant} voir que l'on a une \uncertain{description} \emph{canonique} des \ill{} $(b,c,b',c',t)$ dans une alg. d'Azumaya, \uncertain{pourvu que}, en regardant l'alg.
\[ M(b,c,b',c',t) \]
définie par \uncertain{générateurs} \ill{}, \ill{} d'une base formelle
\[ \underline{1},\ \underline{u},\ \underline{v},\ \underline{uv} \]
et une multiplication explicite. Il \uncertain{resterait} à prouver que celle-ci \uncertain{est} \uncertain{d'}Azumaya. Il suffit de le voir sur les fibres \uncertain{géométriques}, \ill{} on est \uncertain{ramené} au cas d'un corps de base alg. clos. Mais là, l'argument …
\marginal{\ill{} \uncertain{ce qui} \ill{} \uncertain{précède} \ill{}}
\marginal{Attention ! il \uncertain{reste} à vérifier \uncertain{les} conditions \uncertain{d'associativité} \ill{} de $M(b,c,b',c',t)$ \ill{}}
\note{les deux notes marginales sont écrites en biais dans la marge gauche ; on n'en lit que des fragments.}

\page{123}
de \uncertain{fidélisabilité} nous donne en \uncertain{même} temps que $A\simeq M(b,c,b',c',t)$, laquelle est donc bien Azumaya. On gagne !

Mais j'ai oublié de faire la vérification (sur le \struck{alg} corps alg. clos) quand $\delta\delta'=0$, donc $H$ ou $H'$ est \og unipotente\fg{} (\struck{\ill{}} disons $H'$), i.e.
\struck{Considérons le discriminant de ce \ill{} (comme \ill{} sur $k$ alg. clos) $Q(t)$ \uncertain{longuement} comme un polynôme en}
\[ H' = \text{\struck{\ill{}}}\ k[V]/V^2 \]
\marginal{\struck{$\subset A$ s.t. \ill{}} $V$ nilpotent i.e. $\det V=0$, $\operatorname{Tr} V=0$}

La formulation qui suit montre que pour $H$, $H'$ donnés, le \add{problème de} \struck{problème de} \uncertain{fidélisabilité} se \uncertain{dispose} par des choix des \uncertain{éléments} $u\in H$, $v\in H'$ utilisés pour \struck{\ill{}} \uncertain{décrire} \ill{} faire des bases. \struck{On peut donc} Il y a donc trois \og cas types\fg{}
\begin{enumerate}
\item $\delta$ \emph{et} $\delta'$ \struck{\ill{}} $\neq 0$ \qquad OPS $u^2=u$, $v^2=v$ (idempotents) donc $b=b'=1$, $c=c'=0$
\[ Q(t) = t^2 - t = (t-1)t \]
\item $\delta\neq0$, \struck{\ill{}} $\delta'=0$ (ou l'inverse) \qquad OPS $u^2=u$, $v^2=0$, i.e. $b=1$, $c=0$, $b'=0$, $c'=0$
\item $\delta=\delta'=0$ \qquad OPS $u^2=0$, $v^2=0$, i.e. $b=b'=c=c'=0$
\end{enumerate}
Dans les cas 2°) et 3°) et disons quand $v^2=0$, (sans condition sur $u$) on trouve
\[ Q(t) = t^2 \]

\page{124}
Notre assertion de \uncertain{fidélisabilité} \add{\og géom.\fg{} (sur $k$ alg. clos)} revient aux trois assertions suivantes
\begin{itemize}
\item[1°)] Si $t\in k$ \struck{avec $t\notin\{0,1\}$}, alors $\exists$ projecteurs $u$, $v$ \add{(de rang 1)} \struck{en position générale} tels que
\[ \operatorname{Tr} uv = t \]
[et si on suppose $t\notin\{0,1\}$, alors $u$, $v$ \uncertain{seront} en \og position générale\fg{}]
\item[2°)] Si $t\in k$, alors $\exists$ projecteur \add{(rg 1)} $u$, et $v$ nilpotent $\neq0$, \struck{\ill{}} tels que
\[ \operatorname{Tr} uv = t \]
(et si $t\neq0$, on aura $u$, $v$ en pos. gén.)
\item[3°)] Si $t\in k$, alors $\exists$ $u$, $v$ nilp. \struck{\ill{}} $\neq0$, tels que
\[ \operatorname{Tr} uv = t \]
(et si $t\neq0$, $u$, $v$ sont en position générale)
\end{itemize}

Interprétant $A=\operatorname{End}(V)$, $V$ vect. de dim 2, on peut écrire dans tous les cas
\[ u = x'\otimes x , \qquad v = y'\otimes y \qquad\qquad x,y\in V,\ x',y'\in V^{\vee} \]
donc
\[ u\circ v = \langle\text{\struck{$y$}},\text{\struck{$x'$}}\rangle\, \text{\struck{$y'$}}\otimes x \]
\[ t = \operatorname{Tr} uv = \langle x,y'\rangle\langle y,x'\rangle \]
\note{la formule pour $u\circ v$ est surchargée et en partie biffée ; on ne lit pas sûrement les arguments retenus.}
\struck{et il faut voir}

\emph{Cas 1} \struck{\ill{}} $\langle x,x'\rangle=1$, $\langle y,y'\rangle=1$.
On choisit base $e_1,e_2$ de $V$ tel que $x=e_1$, $x'=e_1'$. Si les coeff. de $y$, $y'$ sont $(y_1,y_2)$, $(y_1',y_2')$, la relation $\langle y,y'\rangle=1$

\page{125}
s'écrit
\[ y_1y_1' + y_2y_2' = 1 \]
et on a
\[ t = y_1y_1' \]
\struck{On peut prendre} Les relations deviennent \struck{$y_1=1$, $y_1'=t$} et $y_2$, $y_2'$ tels que $y_2y_2'=1-t$,\note{les indices de cette dernière relation sont mal formés ; la lecture $y_2y_2'$ est imposée par le système qui suit.}
\[ \begin{cases} y_1y_1' = t \\ y_2y_2' = 1-t \end{cases} \]
Si $t\neq0,1$ on peut choisir arbitrairement $y_1,y_2 \in k^*$, et \uncertain{prendre} \struck{$y_1',y_2'$}
\[ y_1' = t/y_1 , \qquad y_2' = \frac{1-t}{y_2} \]
Si $t=0$, \struck{\ill{}} on a
\[ y_1y_1' = 0 , \qquad y_2y_2' = 1 \]
donc familles de solutions \add{: deux \uncertain{p.s.} \uncertain{suivant}} que $y_1=0$ ou $y_1'=0$, et symétriques si $t=1$.

\emph{Cas 2} On peut \uncertain{prendre} $x=e_1$, $x'=e_1'$, relations
\[ y_1y_1' + y_2y_2' = 0 , \qquad t = y_1y_1' \]
i.e. comme \ill{} équations, en les $y_i$, $y_i'$, \struck{\ill{}}
\[ \begin{cases} y_1y_1' = t \\ y_2y_2' = -t \end{cases} \]

\emph{Cas 3} On peut prendre $x=e_1$, $x'=e_2'$.
\struck{\ill{}}
On doit avoir
\[ \begin{cases} y_1y_1' + y_2y_2' = 0 \\ t = y_1'y_2 \end{cases} \]
…

\page{126}
\emph{NB} Finalement, on n'a pas eu besoin du \struck{\ill{}} de faire intervenir $\rho$ dans la vérification, c'était une complication inutile.

Je vais reprendre la théorie précédente. Soient $H$, $H'$ deux fibrés vectoriels de rang 2 \struck{sur} $k$, munis d'une \struck{forme q} structure de $k$-algèbre (assoc. commut.) d'\struck{\ill{}} $H$ et $H'$ formes \emph{\uncertain{tordues}} et tirons
\[ \Delta_H : H\to k , \qquad T_H : H\to k , \qquad \varepsilon_H : k\to H \]
\[ \Delta_{H'} : H'\to k , \qquad T_{H'} : H'\to k , \qquad \varepsilon_{H'} : k\to H' \]
l'unité de ce qui a suivi, \ill{} on n'aura pas besoin de la structure multiplicative, mais seulement des \uncertain{formes} \uncertain{précédentes} \ill{}, \uncertain{seulement} de $T_H$) La formule
\[ \varphi_{\Delta_H}(u,v) = \operatorname{Tr}(u)\operatorname{Tr}(v) - \underbrace{\operatorname{Tr}(uv)}_{\psi_H(u,v)} \]
\[ \varphi_{\Delta_H} = \varphi_H \]
\uncertain{montre} que la forme $\operatorname{Tr}(uv)$ sur $H$ s'exprime par la connaissance de $N_H$ et $T_H$ comme
\[ \operatorname{Tr}_H(uv) = T_H(u)T_H(v) - \varphi_{\Delta_H}(u,v) \]
i.e.
\[ \varphi_H(u,v) + \psi_H(u,v) = T_H(u)T_H(v) \]
\marginal{$\varepsilon_H(\lambda)=\lambda 1_H$, $\varepsilon_{H'}(\lambda) = \rho_0\lambda.1_{H'}$ …{} \emph{NB} $T_H$ se déduit de la donnée de $1_H\in H$ (i.e. $\varepsilon_H$) et de $\Delta_H$, comme $T_H(u)=\varphi_H(u,1)$ \ill{}, \uncertain{du fait} que $\Delta_H(1)=1$ (cf. \ill{} $T(H)=2$) \ill{} \uncertain{hypothèse} \ill{} $\forall u$ \ill{} $\varphi_H(1,u)=\psi_H(1,u)$ \ill{}}
\note{la note marginale est écrite en biais ; des traits la relient au texte. La lecture de $\rho_0$ dans la formule pour $\varepsilon_{H'}$ est douteuse.}

\page{127}
\[ B = H \amalg_k H' = H\oplus H' / k(1_H \ominus 1_{H'}) \]
\note{après la barre de quotient, une première écriture du sous-module ($k(1_H + 1_{H'})$, semble-t-il) est biffée et récrite au-dessus.}
ainsi $H$ et $H'$ se \uncertain{réalisent} comme sous-fibrés vect. de rg 2 de $B$, \ill{}
\[ H\cap H' = k1_B , \qquad H+H' = B \]
\[ 1_B = 1_H = 1_{H'} \quad \text{\uncertain{i.e. image commune}} \]
de $1_H$ et $1_{H'}$ \struck{dans} $B$.

On considère \uncertain{ensuite}
\[ P = H\otimes_k H' \]
et les \uncertain{plongements}

\begin{tikzcd}
 & P & \\
H \arrow[ur, "i_0"] & & H' \arrow[ul, "j_0"']
\end{tikzcd}

\[ i(u) = u\otimes 1 , \qquad j(v) = 1\otimes v \]
On peut \uncertain{considérer} $B$ comme un sous-fibré vectoriel de $P$, \struck{\ill{}} [i.e. \struck{\ill{}} \uncertain{conoyau} \ill{} $k$
\[ i(\lambda 1_H) = j(\lambda 1_{H'}) = \lambda\, \underbrace{1\otimes 1}_{\overset{\text{déf}}{=}\,1_P} \]
] \uncertain{de} \uncertain{façon} \uncertain{unique} \ill{} $B\to P$

\begin{tikzcd}
 & P & \\
 & B \arrow[u] & \\
H \arrow[ur] \arrow[uur, "i_0"] & & H' \arrow[ul] \arrow[uul, "j_0"']
\end{tikzcd}

\emph{NB} $B \xrightarrow{T_B} k$ \uncertain{induisant} $T_H$, $T_{H'}$ \marginal{cf. \ill{} $T_{H'}(1)=2$}
\marginal{Mais \ill{} \ill{} $\varepsilon_H$, $\varepsilon_{H'}$ \ill{} $\Delta_H(1)=\Delta_{H'}(1)=1$ \ill{}}
\marginal{Attention ! \ill{} \uncertain{structure} \ill{} \uncertain{d'algèbre} \ill{}}
\note{les deux premières notes marginales sont écrites en biais dans la marge gauche, la troisième en bas à gauche.}

\page{129}
Supposons maintenant que \struck{$H$, $H'$} l'on se donne une alg. Azumaya $A$ et des plongements

\begin{tikzcd}
 & A & \\
H \arrow[ur, "i"] & & H' \arrow[ul, "j"']
\end{tikzcd}

\marginal{plus généralement, $A$ dans une \uncertain{catégorie} \uncertain{arbitraire} \uncertain{fournissant} H.C. d'\ill{} $\Delta_A$, $\operatorname{Tr}_A$ …}
On peut alors définir \uncertain{un} \uncertain{homomorphisme} \uncertain{de} \uncertain{modules}
\[ P \xrightarrow{\ f\ } A \]
par la condition
\[ f(u\otimes v) = i(u).j(v) \]
\struck{ce qui induit et la donnée de} et on récupère $i$, $j$ par
\[ i = f\circ i_0 , \qquad j = f\circ j_0 \]
La donnée de $i$, $j$ équivaut à la donnée de $f$, soumis à la condition que $f|H$ et $f|H'$ \uncertain{soient} des hom. d'algèbres. On a alors

\begin{tikzcd}
 & A & \\
 & P \arrow[u, "f"] & \\
 & B \arrow[u] & \\
H \arrow[ur] & & H' \arrow[ul] \\
 & k \arrow[ul, "\varepsilon_H"] \arrow[ur, "\varepsilon_{H'}"'] &
\end{tikzcd}

\marginal{Attention, on ne suppose pas que $f$ soit un iso de Mod, i.e. que $H$, $H'$ engendrent $A$ comme algèbre}

\page{130}
Cette situation permet de définir sur $P$
\begin{itemize}
\item forme linéaire $L_P = \operatorname{Tr}_A\circ f$, qui prolonge $T_B$
\item forme quadratique $\Delta_P = \det_A\circ f$
\item forme bilinéaire sym. $\psi_P$ polarisée de $\Delta_P$, $\psi_P = \psi_A\circ f$ (où $\psi_A$ est polarisée de $\det_A$)
\item forme bilinéaire $\varphi_P = \varphi_A\circ f$
\end{itemize}
i.e.
\[ \varphi_P(U,V) = \operatorname{Tr}(f(U).f(V)) \]
\marginal{\emph{NB} On a \struck{\ill{}}
$\varphi_P(U,V)+\psi_P(U,V) = L_P(U)L_P(V)$ et $\varphi_P(1,U)=\psi_P(1,U)=L_P(U)$}

D'autre part, on peut restreindre ceci à $B$, \uncertain{d'où}
forme quadr. $\Delta_B$, forme bil. sym. $\psi_B$, $\varphi_B$.

En restreignant encore ces trois formes sur $B$ à $H$, on trouve les trois formes dont $H$ est déjà équipée, $\Delta_H=\det_H$, $\psi_H$ = forme polarisée de $\Delta_H$, $\varphi_H = ((u,v)\mapsto T_H(uv))$, et itou pour $H'$.

On a d'autre part
\[ (1) \qquad \Delta_P(u\otimes v) = \Delta_H(u)\,\Delta_{H'}(v) \]
On se donne $H$, $H'$ comme ci-dessus.

\emph{Proposition} (Considérons les \add{quatre} \struck{\ill{}} ensembles suivants :
\begin{itemize}
\item[$\underline{Q}_P$ :] \struck{\ill{}} ens. des formes quadratiques $\Delta_P$ sur $P$ satisfaisant (1) ci-dessus
\item[$\underline{Q}_B$ :] \struck{\ill{}} ens. des formes quadr. $\Delta_B$ sur $B$ prolongeant $\Delta_H$ et $\Delta_{H'}$
\item[$\underline{\Psi}_B$ :] \struck{\ill{}} ens. des formes bil. sym. $\psi_B/\Delta_B$ sur $B$ qui prolongent $\psi_H$, $\psi_{H'}$
\item[$\underline{\Phi}_B$ :] \struck{\ill{}} ens. des formes bil. symétriques $\varphi_B$ sur $B$
\end{itemize}
\note{la liste se poursuit en haut de la p.~132.}

\page{131}
2) \struck{On voit que} \uncertain{Posons}
\[ L = H/k \simeq B/H' , \qquad L' = H'/k \simeq B/H \]
\struck{et considérons une forme}
\[ \text{\struck{$\lambda : L\otimes L' \to k$}} \]
\struck{Si $\Delta_B\in\underline{Q}_B$, on peut définir}
\note{ces trois lignes sont barrées d'un trait ondulé.}
On a
\[ (L\otimes L')^{\vee} \simeq \text{Esp. des formes bil. \add{sym.} sur $B$ nulles sur $H'$ \struck{\ill{}} et $H$ \struck{\ill{}}} \]
i.e. qui passent au \uncertain{quotient} \struck{\ill{}} $B/k\simeq L\oplus L'$ et qui sont nulles sur $L$ et sur $L'$, donc de la forme
\[ \bigl((x,y),(\xi,\eta)\bigr) \longmapsto \lambda(x\otimes\eta) + \lambda(\xi\otimes y) \qquad x,\xi\in L,\ y,\eta\in L' \]
avec $\lambda : L\otimes L'\to k$.
\marginal{car $B = H+H'$}

Donc $\underline{\Psi}_B$ est un torseur sous $(L\otimes L')^{\vee}$, et de même $\underline{\Phi}_B$. Itou pour $\Delta_B$, car
\[ (L\otimes L')^{\vee} \simeq \text{formes quadratiques sur $B$ qui sont nulles sur $H$, $H'$} \]
(donc \uncertain{passent} au quotient \ill{} $B/k\simeq L\oplus L'$ \uncertain{en} des formes quadratiques nulles \struck{sur $L$, $L'$, donc de la forme} \ill{}
$\Delta(x\oplus y) = \lambda(x\otimes y)$).
Les applications $\beta$, $\gamma$ sont compatibles avec ces structures de torseurs. Cela suffit : comme \uncertain{quelles} sont iso.

\page{132}
\begin{itemize}
\item[] qui prolongent $\varphi_H$, $\varphi_{H'}$ ;
\item[$\underline{L}_P$ :] Ens. des formes \emph{lin.} sur $P$ qui prolongent $T_B$
\end{itemize}
On a des flèches canoniques
\note{au-dessus du diagramme, une première version (avec $\underline{\Phi}_P$, $\underline{L}_P$, $\underline{\Phi}_B$) est raturée.}

\begin{tikzcd}
\underline{Q}_P \arrow[r, "\alpha"] \arrow[dr, "\delta"'] & \underline{Q}_B \arrow[r, "\beta"] & \underline{\Psi}_B \arrow[r, "\gamma"] & \underline{\Phi}_B \arrow[dll] \\
 & \underline{L}_P & &
\end{tikzcd}

que voici : $\alpha$ est la restriction $\alpha(\Delta_P) = \Delta_P\circ\alpha$,\note{lecture de la formule douteuse : il faut sans doute entendre la restriction de $\Delta_P$ à $B$.} $\beta$ associe à la forme $\Delta_B\in\underline{Q}_B$ la forme bil. sym. associée $\psi_B$, on définit $\gamma$ par la formule
\[ \underbrace{\gamma(\psi)}_{\overset{\text{déf}}{=}\ \varphi_B} + \psi_B = \bigl((u,v)\longmapsto T_B(u)\,T_B(v)\bigr) \]
on définit $\delta$ par
\[ \delta(\Delta_P) = \{ u \longmapsto \varphi_{\Delta_P}(u,1_P) \} \]
Ceci dit, les quatre applications $\alpha$, $\beta$, $\gamma$, $\delta$ sont \emph{bijectives}.

\emph{Remarques} 1) Pour définir les cinq ensembles $\underline{Q}_P$, $\underline{L}_P$, $\underline{Q}_B$, $\underline{\Psi}_B$, $\underline{\Phi}_B$, et les applications $\alpha$, $\beta$, $\gamma$, $\delta$, on n'a utilisé sur $H$, $H'$ que la donnée de $\Delta_H$, $\Delta_{H'}$ et $1_H$, $1_{H'}$ ($\to T_H$, $T_{H'}$, cf. \ill{}) satisfaisant
\[ \Delta_H(1_H) = \Delta_{H'}(1_{H'}) = 1 . \]
\marginal{On définit $\zeta$ : \ill{} ainsi : $\underline{L}_P$ s'identifie \uncertain{aux} formes \emph{bilin.} sur $H\times H'$ \uncertain{telles} que leurs restrictions à $1\times H'$ et à $H\times1$ soient données par $T_{H'}$ et $T_H$. \uncertain{À} $\varphi_B$ on associe $\zeta(\varphi_B) : (u,v)\mapsto\varphi_B(u,v)$. C'est \uncertain{encore} un \uncertain{iso} de torseurs sous $(L\otimes L')^{\vee}$, donc iso.}

\page{133}
\note{le haut de la page achève l'énoncé du lemme qui termine la p.~134 (« … $\to \underline{\mathrm{Hom}}(\det\check{H},\det H')$ ») ; la p.~134 semble donc se lire entre les pp.~132 et 133, et la fin de la présente page (« un sixième ensemble ») se poursuit en p.~135.}
$\simeq (\det H)\otimes\det(H')$. Cette \struck{\ill{}} fonction quadratique est \emph{universelle}, pour les fonctions quadratiques
\[ P \xrightarrow{\ q\ } Q \]
telles que
\[ q(x\otimes y) = 0 \qquad \forall x,y \in P . \]

Ceci nous montre (\struck{que, \ill{}} \uncertain{grâce} à \struck{$\underline{Q}_P$} la situation précédente) que si $\underline{Q}_P$ n'est pas vide, c'est un torseur sous \struck{\ill{}} $(L\otimes L')^{\vee}$. Il faut vérifier que $\alpha$, $\delta$ sont compatibles aux structures torseurs --- moyennant quoi on voit que c'est des iso.

\struck{Pour vérifier que $\underline{Q}_P$ n'est pas vide, on prend sur $P$ la structure \struck{tensorielle} d'alg. produit de celles de $H$, $H'$, \struck{et la forme} \uncertain{donnée} de H.C., dont \ill{} la fonction dét fait l'affaire.}
\note{ce paragraphe est barré de traits croisés.}
On peut utiliser la construction des algèbres $M(b,c,b',c',t)$. Cela nous fait songer à un sixième ensemble

\page{134}
en fait, comme on a
\[ H\to L , \qquad H'\to L' \]
on en déduit
\[ P \overset{\text{déf}}{=} H\otimes H' \longrightarrow L\otimes L' \quad \text{\uncertain{épi}} \]
nul sur $1\otimes H'$ et sur $H\otimes 1$, donc sur $B$, d'où une suite exacte
\[ 0 \to B \to P \to L\otimes L' \to 0 \]
\struck{Ceci dit} Ainsi les formes linéaires \add{$T_P$} sur $P$ dont la restr. à $B$ est $T_B$, forment \add{$\underline{L}_P$} un torseur sous $(L\otimes L')^{\vee}$. L'application
\[ \zeta : \underline{\Phi}_B \Longrightarrow \underline{L}_P \]
\[ \varphi_B \longmapsto \bigl( x\otimes y \mapsto \varphi_B(x,y) \bigr) \]
est compatible avec les opérations de $(L\otimes L')^{\vee}$. Donc $\zeta$ est iso.
\marginal{i.e. les formes bil. $H\times H'\to k$ \uncertain{telles} que $f(u,1)=T_H(u)$, $f(1,v)=T_{H'}(v)$}

Notons le \emph{Lemme} Soient $H$, $H'$ deux modules \uncertain{localement} \uncertain{libres} de rang 2 sur $k$, $P\simeq H\otimes H'$. On définit une fonction quadratique canonique
\[ P \xrightarrow{\ \Delta_{H,H'}\ } \det H \otimes \det H' \]
en interprétant
\[ P \simeq \underline{\mathrm{Hom}}(\check{H},H') \to \underline{\mathrm{Hom}}(\det\check{H},\det H') \]

\page{135}
$AZ_P$ = structures de $k$-algèbre de H.C. sur $P=H\otimes H'$, induisant sur $1\otimes H'$ et $H\otimes1$ les structures d'alg. données.

\begin{tikzcd}[column sep=small, row sep=large, nodes={font=\scriptsize}]
AZ_P \arrow[r] \arrow[drr, "\alpha\mapsto\operatorname{Tr}_\alpha"'] & \underline{Q}_P \arrow[r] \arrow[dr, dashed] & \underline{Q}_B \arrow[r] & \underline{\Psi}_B \arrow[r] & \underline{\Phi}_B \arrow[dll, "{\varphi\mapsto(x\otimes y\mapsto\varphi(x,y))}"] \\
 & & \underline{L}_P & &
\end{tikzcd}

\note{étiquettes des flèches de la ligne du haut, de gauche à droite : $\alpha\mapsto\det_\alpha$ ; $\Delta_P\mapsto\Delta_P|B$ ; $\Delta_B\mapsto\psi_{\Delta_B}$ ; $\varphi\mapsto\psi$ t.q. $\varphi+\psi=T_B\otimes T_B$. La flèche en tirets $\underline{Q}_P\to\underline{L}_P$ porte $\Delta_P\mapsto(U\mapsto\psi_{\Delta_P}(1,U))$, avec un mot illisible ; dans la dernière étiquette, $x\in H$, $y\in H'$.}

La commutativité du triangle total \uncertain{est} \uncertain{immédiate}, ainsi que celle des \ill{} \uncertain{entraîne} \uncertain{celle} du triangle de \uncertain{gauche}, \ill{} \uncertain{donc} celle du triangle de droite (qui n'était pas \uncertain{évidente} \uncertain{ainsi}).

\struck{L'on} On aurait pu aussi introduire $\underline{\Psi}_P$ et $\underline{\Phi}_P$ (on les définirait \uncertain{comme} \uncertain{ci-dessus}) mais passons ! Parmi les 6 \struck{\ill{}} ens., celui qui est de \uncertain{description} alg. la plus simple est
\[ \underline{L}_P \subset \mathrm{Bil}(H,H';k) . \]
On \uncertain{va} \uncertain{dire} $T_P\in\underline{L}_P$ un \uncertain{twist} \struck{\ill{}} \uncertain{des} algèbres \struck{\ill{}} quadratiques $H$, $H'$. L'objet le plus \struck{fondamental} \uncertain{toutefois} \ill{} \uncertain{est} la structure d'alg. de H.C. sur $P = H\otimes H'$,
\marginal{C'est défini en termes de $H$, $H'$ munis de $T_H$, $T_{H'}$ et de $1_H$, $1_{H'}$, \ill{} (à l'exclusion de la donnée de $\Delta_H$, $\Delta_{H'}$)}

\page{136}
qui donne \uncertain{aussitôt} les \ill{} \uncertain{structures}
\[ \Delta_P,\ T_P,\ \varphi_P,\ \psi_P,\ \Delta_B,\ \psi_B,\ \varphi_B \]

\emph{Th. récapitulatif} ① Les flèches du diagramme de la p. précédente sont bijectives, le diagramme est commutatif, et les \add{5} flèches qui en sont \uncertain{issues} \uncertain{de la} source $AZ_P$ sont compatibles avec les structures de torseur sous $(L\otimes L')^{\vee}$.

② Pour qu'une structure \add{d'alg. sur $P$} associée à un twist en fasse une alg. d'Azumaya, il faut et il suffit que la forme $\Delta_B$ sur $B$ soit non dégénérée, i.e. que son discriminant
\[ \operatorname{disc}'(\Delta_B) \in (\det B)^{\otimes-2} = (L\otimes L')^{\otimes-2} \]
soit inversible.

③ \uncertain{L'application}
\[ \underline{L}_P \xrightarrow{\ \delta_{P,B}\ } (L\otimes L')^{\otimes-2} \]
\note{sous la flèche : « discriminant divisé de la forme quad. sur $B$ associée à un twist ».}
est une application \struck{quadratique} pol. de degré 2
\marginal{\ill{} \uncertain{cette} \ill{} \uncertain{définition} \ill{} $1_H$, $1_{H'}$, $T_H$, $T_{H'}$ \ill{} \uncertain{applications}, mais $\delta$ \uncertain{fait} \uncertain{appel} à $\Delta_H$, $\Delta_{H'}$ ?}

\page{137}
dont la forme quadratique associée
\[ \underbrace{(L\otimes L')^{-1}}_{\text{esp. des translations de } \underline{L}_P} \longrightarrow (L\otimes L')^{-2} \]
\note{au-dessus de la flèche, un « $\delta_{P,B}$ » est raturé.}
est l'élévation au carré (donc $\delta_{P,B}$ est \struck{après passage à} un morphisme fini surjectif de degré 2).

Prenons l'image inverse de la section nulle, c'est un schéma quadratique \struck{fini} sur $S=\operatorname{Spec} k$ (\uncertain{donc} \uncertain{défini} \uncertain{par} \add{qui est défini par une alg. \ill{} quadratique} $k'$), cela dit, il y a un morphisme canonique \struck{\ill{}}
\[ \boxed{\, X\times_S X' \xrightarrow{\ \alpha\ } \tilde{S} \,} \]
\struck{tel que} tel que
\[ \alpha(\bar{x},x') = \alpha(x,\bar{x}') = \overline{\alpha(x,x')} \]
(où $x\mapsto\bar{x}$ etc. désignent les inv. canoniques) d'où \uncertain{pour} $X$, $X'$ étales \uncertain{(car)} $\tilde{S}$ est \uncertain{étale} d'\ill{} \add{produit contracté des $(\mathbb{Z}/2\mathbb{Z})_S$-torseurs}
\[ X\wedge_S X' \longrightarrow S , \]
\struck{défini} hom. de $(\mathbb{Z}/2\mathbb{Z})_S$-torseurs, donc iso.
\marginal{$X$, $X'$ \uncertain{des} $S$ \uncertain{définis} \uncertain{par} $H$, $H'$ ; \uncertain{posé} $\tilde{S} = X\wedge X'$}

Il faut définir l'application $\alpha$. Soient
\[ \varphi : H\to k , \qquad \varphi' : H'\to k \]
\struck{deux hom. d'alg., considérons la forme}
\[ \text{\struck{$\varphi\otimes\varphi' : H\otimes H' = P\to k \qquad u\otimes v\mapsto \varphi(u)\varphi'(v)$}} \]

\page{138}
Choisissons bases $(1,u)$ dans $H$, $(1,v)$ dans $H'$ (d'où $b,c,b',c'$), $\varphi$ et $\varphi'$ sont données par
\[ \varphi(u) = z , \qquad \varphi'(v) = z' \]
avec $z\in k$
\[ z^2 - bz + c = 0 \qquad z'^2 - b'z' + c' = 0 , \]
et on cherche $t$ solution de
\[ t^2 - bb't + (b^2b' + b'^2b - 4cc') = 0 . \]
\note{sic : le terme constant est écrit $b^2b'+b'^2b-4cc'$ ; la ligne suivante et la p.~139 portent $b^2c'+b'^2c-4cc'$.}
Mais le discriminant de cette équation est
\[ \Delta = (bb')^2 - 4(b^2c' + b'^2c - 4cc') = \underbrace{(b^2-4c)}_{\delta}\,\underbrace{(b'^2-4c')}_{\delta'} \]
Or \struck{\ill{}} les \uncertain{formules} \struck{\ill{}} \uncertain{heuristiques} (\uncertain{valables} sur un corps de car $\neq2$)
\[ 2z = b\pm\sqrt{\delta} , \qquad 2z' = b'\pm\sqrt{\delta'} \]
\struck{donc}
\[ \delta = (2z-b)^2 \qquad \delta' = (2z'-b)^2 \]
\note{sic $(2z'-b)^2$ pour $(2z'-b')^2$.}
(\uncertain{vrai} \uncertain{sans} condition \uncertain{corps} ou car. \ill{}) et les solutions de l'équation en $t$ \add{devraient} \uncertain{satisfaire} (sur corps car $\neq2$)
\[ \tfrac14 (2t-bb')^2 = \Delta = \delta\delta' = \bigl((2z-b)(2z'-b')\bigr)^2 \]
donc on peut \struck{poser} supposer
\[ 2t - bb' = \pm(2z-b)(2z'-b') \]
soit
\[ 2t = bb' \pm (2z-b)(2z'-b') . \]
Choisissant \struck{\ill{}} le $\pm$, on trouve
\[ 2t = -4zz' + 2(b'z+bz') \]
soit
\[ \boxed{\, t = (b'z+bz') - 2zz' \,} \]
l'autre racine \add{sera}
\[ t' = \bigl(bb' - (b'z+bz')\bigr) + 2zz' \qquad (= bb'-t) \]

\page{139}
Il suffit de vérifier que l'on a bien
\[ tt' = b^2c' + b'^2c - 4cc' \]
(ce qui \uncertain{exige} \ill{} $t+t'=bb'$). Vérification \uncertain{immédiate} (il \uncertain{n'est} pas besoin de \uncertain{la} \uncertain{forme}).

On peut l'écrire
\[ \boxed{\, t_{\varphi,\varphi'}(u,v) = \operatorname{Tr}(v)\varphi(u) + \operatorname{Tr}(u)\varphi'(v) - 2\varphi(u)\varphi'(v) \,} \]
Or, on a là une forme bil. en $u$, $v$, qui pour $v=1$ se réduit :
\[ t(u,1) = 2\varphi(u) + \operatorname{Tr}(u) - 2\varphi(u) = \operatorname{Tr}(u) \]
et de même pour $u=1$ \struck{\ill{}}
\[ t(u,1) = \operatorname{Tr} u , \qquad t(1,v) = \operatorname{Tr} v \]
C'est donc un élément de $\underline{L}_P$ ! Soit
\[ t_{\varphi,\varphi'} \in \underline{L}_P \]
\struck{Je dis que} défini par $\varphi$, $\varphi'$. Je dis que
\[ \delta(t_{\varphi,\varphi'}) = 0 \]
Pour le voir, on peut faire le calcul \uncertain{avec} une base \struck{\ill{}} $(1,u)$, $(1,v)$ telle que
\[ \varphi(u) = \varphi(v) = 0 \]
d'où $\det u = \det v = 0$ \add{ou $c=c'=0$} \struck{\ill{}} de $z^2-bz+c=0$, de sorte que maintenant
\[ t(u,v) = 0 . \]
\struck{\ill{}} On va a l'identité \struck{\ill{}}
\[ \delta(u,v,t) = t^2 - bb't + (b^2c' + b'^2c - 4cc') = t(t-bb') \quad \text{car } c=c'=0 \]
\struck{\ill{}} donc en y \struck{faisant $t$} $t(u,v)=0$ on \uncertain{a} bien une solution.

\page{140}
\struck{Alors} On \uncertain{peut} bien définir ainsi
\[ X\times X' \longrightarrow X\wedge X' \qquad (\overset{\text{déf}}{=} \tilde{S}) \]
(car la définition précédente commute au chgt. de base).

Notons qu'on a dans $\underline{L}_P$ une involution canonique,
\[ T \longmapsto T' \]
où
\[ T'(u,v) = T(u',v) = T(u,v') \]
(Il faut prouver que l'on a bien
\[ \boxed{\, T(u',v) = T(u,v') \,} \overset{\text{déf}}{=} T'(u,v) \]
pour \struck{\ill{}} forme $T\in\underline{L}_P$, \struck{alors on} \uncertain{tenons} des plongements $H,H'\subset A_T$, cela signifie simplement
\[ \operatorname{Tr}_{A_T}(u'v) = \operatorname{Tr}_{A_T}(uv') \]
or le premier membre, compte tenu que $\operatorname{Tr}_{A_T}$ est invariant sous $'$, est aussi $\operatorname{Tr}_{A_T}((u'v)'=v'u)$ et $\operatorname{Tr} v'u = \operatorname{Tr} uv'$ OK.

On a donc
\[ (T')' = T , \]
et on voit aussitôt que
\note{la phrase se poursuit au-delà de la fin du lot.}

\end{document}
